\documentclass[11pt]{article}

\usepackage[a4paper,margin=1in]{geometry}
\usepackage{amsmath,amssymb,amsthm,mathtools}
\usepackage{booktabs}
\usepackage{enumitem}
\usepackage{microtype}
\usepackage{xcolor}
\usepackage{hyperref}

\hypersetup{
    colorlinks=true,
    linkcolor=blue!60!black,
    citecolor=blue!60!black,
    urlcolor=blue!60!black,
    pdftitle={Shadow Trigonometry},
    pdfauthor={Vinicius Santos}
}

\newtheorem{definition}{Definition}[section]
\newtheorem{theorem}[definition]{Theorem}
\newtheorem{corollary}[definition]{Corollary}
\newtheorem{proposition}[definition]{Proposition}
\newtheorem{lemma}[definition]{Lemma}
\newtheorem{observation}[definition]{Observation}
\newtheorem{question}[definition]{Question}
\newtheorem{remark}[definition]{Remark}

\newcommand{\B}{\mathsf{B}}
\newcommand{\Tr}{\operatorname{Tr}}
\newcommand{\Nm}{\operatorname{N}}
\newcommand{\Sp}{\operatorname{Sp}}
\newcommand{\Phase}{\operatorname{PhaseLift}}
\newcommand{\Fix}{\operatorname{Fix}}
\newcommand{\ph}{\varphi}
\newcommand{\ps}{\psi}
\newcommand{\sig}{\sigma}
\newcommand{\Z}{\mathbb{Z}}
\newcommand{\Q}{\mathbb{Q}}
\newcommand{\R}{\mathbb{R}}
\newcommand{\C}{\mathbb{C}}
\newcommand{\T}{\mathbb{T}}

\title{Shadow Trigonometry\\
\large Norm-One Orbits, Golden Traces, and Cyclotomic Phase}
\author{Vinicius Santos\\
\small Exploratory draft}
\date{June 2026}

\begin{document}
\maketitle

\begin{abstract}
The companion manuscripts \emph{A Golden Shadow Algebra} and \emph{Quotients Without Division} suggest a recurring principle: operations that look primitive in ordinary analytic notation may become projections, relations, or constraints in a $\ph$-native algebraic presentation.  Addition becomes trace projection; quotient-like operations split into unit action, fixed averaging, approximation, relation solving, and factor projection.  This note applies the same strategy to trigonometry.

Classical trigonometry usually begins with angles, sine, cosine, tangent, and the constant $\pi$.  We instead begin with a norm-one phase object $U$ satisfying $U^\sig=U^{-1}$ and study its unnormalized trace and spread shadows
\[
    T(U)=U+U^{-1},
    \qquad
    S(U)=U-U^{-1}.
\]
The main thesis is: phase first, trace second, angle later.  The algebraic core of trigonometry consists of norm-one constraints, torsion equations, unit orbits, and trace/spread projections.  Classical sine, cosine, tangent, angle, and $\pi$ appear later as normalized or analytic coordinatizations.

The central result is the golden circular bridge.  If a norm-one phase $U$ satisfies
\[
    U+U^{-1}=\ph,
\]
then the trace recurrence, equivalently the tenth cyclotomic polynomial, gives
\[
    U^5=-1,
    \qquad
    U^{10}=1,
\]
and in the integral-domain setting the phase has exact order $10$.  Thus the golden ratio is not only a real quadratic unit; it is also the visible trace of a hidden pentagonal circular phase.  In classical language this is the identity $2\cos(\pi/5)=\ph$, but the shadow formulation reverses the emphasis: $\ph$ is the scalar shadow, and angle is a later coordinate on the phase lift.
\end{abstract}

\tableofcontents

\section{Introduction: trigonometry without starting from angles}

The broader sequence was motivated by Odrzywo\l{}ek's EML construction, where the leaf $1$ and a single binary analytic operator generate the elementary scientific-calculator repertoire~\cite{Odrzywolek2026EML}.  The first companion paper, \emph{A Golden Shadow Algebra}, then asked what changes when the leaf is the golden ratio rather than $1$, and began with a simple scalar operation
\[
    \B(x,y)=xy-y=y(x-1)
\]
and the golden leaf
\[
    \ph=\frac{1+\sqrt5}{2}.
\]
It showed that the scalar grammar generated by $\B$ and $\ph$ is multiplication/predecessor-native but not addition-native.  Addition appears only after passing to a typed pair layer:
\[
    x+y=\Tr\langle x\mid y\rangle.
\]

The second paper, \emph{Quotients Without Division}, applied the same discipline to division.  Once the exact scalar constants are known to be
\[
    \Z[\ph]
\]
and therefore dense in a chosen real embedding of $\R$ (while discrete under the full Minkowski embedding), division is no longer needed merely to obtain dense constants.  Instead, division splits into roles: unit action, fixed-denominator rate localization, solved averaging, rational approximation, reciprocal approximation, quotient relations, factor projection, and normalization constraints.

The present note turns to trigonometry.  Classical trigonometry bundles together several different notions:
\[
    \sin,\quad \cos,\quad \tan,\quad e^{i\theta},\quad \theta,
    \quad \pi,\quad \text{rotation},\quad \text{periodicity}.
\]
The guiding question is:
\begin{quote}
What remains of trigonometry if we do not begin with angles, normalized sine/cosine, or division-based tangent?
\end{quote}

The answer proposed here is:
\begin{quote}
Trigonometry is projection theory of norm-one shadow orbits.
\end{quote}

The primitive object is not an angle $\theta$.  It is a phase unit $U$ equipped with an involution
\[
    U^\sig=U^{-1}
\]
and a norm-one condition
\[
    \Nm(U)=UU^\sig=1.
\]
The basic projections are unnormalized:
\[
    T(U)=U+U^{-1},
\]
\[
    S(U)=U-U^{-1}.
\]
Only after choosing a complex embedding and dividing by $2$ and $2i$ do these become the classical coordinate functions
\[
    \cos\theta=\frac{T(U)}2,
    \qquad
    \sin\theta=\frac{S(U)}{2i}.
\]
Those normalized coordinates are valuable, but they are not the algebraic starting point.  In the terminology of the quotient manuscript, the divisions by the fixed constants $2$ and $2i$ belong first to a rate/localization layer and, once evaluated, to an averaging/coordinatization layer.  They do not introduce variable-denominator singularities, but they do enlarge the coordinate system beyond the unnormalized trace/spread shadows.

The central slogan is therefore:
\[
\boxed{\text{phase first, trace second, angle later}.}
\]

From this point onward, ``native'' means native to the phase/trace layer, not native to the finite scalar $\B$-grammar alone.  Expressions such as $U+U^{-1}$ and $U-U^{-1}$ are trace/spread projections from a richer phase object, just as scalar addition in the first companion manuscript arose from a typed trace layer.

\subsection*{Main results and claims}

The paper develops the following theorem package.
\begin{enumerate}[label=(\roman*)]
    \item Norm-one phase objects give unnormalized trace and spread projections.
    \item Chord norm satisfies the coordinate-free identity
    \[
        \Nm(U-1)=2-T(U).
    \]
    \item Trace orbits
    \[
        T_n(U)=U^n+U^{-n}
    \]
    satisfy the Chebyshev/Lucas recurrence
    \[
        T_{n+1}=T_1T_n-T_{n-1}.
    \]
    \item A phase lift of a scalar $t$ is the relation
    \[
        U^2-tU+1=0,
    \]
    replacing the angle-first expression $e^{i\arccos(t/2)}$.
    \item For $t=\ph$, the lifted phase satisfies the tenth cyclotomic equation
    \[
        U^4-U^3+U^2-U+1=0,
    \]
    hence
    \[
        U^5=-1,
        \quad U^{10}=1,
    \]
    and, in the integral-domain setting used below, $U$ has exact order $10$.
\end{enumerate}

The last point is the golden circular bridge.  Classically one writes
\[
    2\cos(\pi/5)=\ph.
\]
Here the emphasis is reversed: $\ph$ is a visible scalar, and the hidden phase $U$ satisfying $U+U^{-1}=\ph$ is its circular lift.

The payoff of the reordering is not a new trigonometric identity.  It is a change in where singularity-prone and normalization-dependent operations enter.  Trace, spread, torsion, chord norm, and phase lift are algebraic and unnormalized; angle, $\pi$, tangent quotients, and normalized sine/cosine are delayed to coordinate or analytic completion layers.  This makes the boundary between algebraic phase structure and analytic trigonometric coordinates explicit.

\section{Inherited machinery from the companion manuscripts}

Only a small part of the previous two papers is needed here.

\subsection{Scalar core and dense constants}

The golden scalar operation is
\[
    \B(x,y)=xy-y.
\]
With leaf $\ph$ it gives
\[
    1=\B(\ph,\ph),
\]
\[
    0=\B(1,\ph),
\]
\[
    -x=\B(0,x),
\]
\[
    x-1=\B(x,1),
\]
\[
    xy=\B(x+1,y).
\]
Here $x+1$ is not primitive; it is constructed from predecessor and negation in the scalar grammar.  The exact scalar constants are
\[
    \Z[\ph]=\{a+b\ph:a,b\in\Z\},
\]
and this set is dense in the chosen real embedding of $\R$, while discrete under the full Minkowski embedding into $\R^2$.  Thus the $\ph$-native scalar substrate obtains one-coordinate density without introducing division.

\subsection{Addition as trace projection}

Scalar $\B$-terms do not express mixed addition such as $x+y$.  Pair formation and trace provide the natural additive completion:
\[
    x+y=\Tr\langle x\mid y\rangle.
\]
This is the first instance of the recurring pattern:
\[
    \text{classical operation} \quad\leadsto\quad \text{projection from a richer object}.
\]

\subsection{Quotients as relations and factor projections}

The second companion manuscript decomposes division.  A quotient
\[
    q=\frac uv
\]
may instead be represented by the relation
\[
    u=qv.
\]
Fixed integer division may be carried as an unevaluated rate/localization class and, when evaluated, represented by an averaging section of a diagonal trace map; variable quotient-like expressions may instead be represented by relations or projections.  A derivative-like quotient may be represented by factor projection:
\[
    \Delta f=G\,\Delta x.
\]
The present paper applies the same principle to trigonometry.  Sine, cosine, tangent, and angle are not primitives.  They arise from norm-one phase objects, trace/spread projections, relations, and completions.

\section{The jobs of trigonometry}

Classical trigonometry is not a single operation.  It is a bundle of roles.

\begin{center}
\begin{tabular}{p{0.28\linewidth}p{0.27\linewidth}p{0.35\linewidth}}
\toprule
Classical role & Classical object & Shadow-native replacement \\
\midrule
Unit circle & $x^2+y^2=1$ & norm-one condition $\Nm(U)=1$ \\
Rotation & $e^{i\theta}$ & phase unit $U$ \\
Period & $2\pi$ & torsion $U^n=1$ \\
Angle & $\theta$ & orbit index or phase coordinate \\
Cosine & real projection & trace $U+U^{-1}$ \\
Sine & orthogonal projection & spread $U-U^{-1}$ \\
Tangent & $\sin/\cos$ & spread--trace slope relation \\
Addition formulas & $\sin(a+b),\cos(a+b)$ & multiplication $U^mU^n=U^{m+n}$ \\
Chebyshev recurrences & $T_n(\cos\theta)$ & trace recurrences of powers \\
$\pi$ & analytic period & coordinate of phase completion \\
\bottomrule
\end{tabular}
\end{center}

This table is the organizing map.  The rest of the paper rebuilds the algebraic core of these roles without beginning from normalized sine, cosine, tangent, or angle.

\section{Shadow involutions and norm-one phase objects}

Before defining phase objects, we must clarify the role of the involution $\sig$.

\begin{remark}[The meaning of $\sig$]
In the first companion manuscript, $\sig$ denoted the nontrivial Galois conjugation of the real quadratic field $\Q(\sqrt5)$:
\[
    \ph^\sig=\ps=1-\ph=-\ph^{-1}.
\]
In the present paper, $\sig$ denotes a formal shadow involution appropriate to the phase object under study.  In circular phase examples it acts by inverse reflection:
\[
    U^\sig=U^{-1}.
\]
The common structure is not a single ambient field automorphism reused in every setting.  It is an involutive shadow operation whose projections define trace, norm, and spread.  In both cases, however, the Vieta pattern is the same: the involution swaps the two roots of the defining quadratic.  For the real golden pair the roots satisfy $r_1+r_2=1$ and $r_1r_2=-1$; for a circular phase lift they satisfy $U+U^{-1}=t$ and $UU^{-1}=1$.
\end{remark}

This distinction matters.  In the real golden pair,
\[
    \ph^\sig=\ps\ne \ph^{-1}.
\]
Indeed,
\[
    \ps=-\ph^{-1}.
\]
In a circular norm-one phase, however,
\[
    U^\sig=U^{-1}.
\]
The notation is shared because the projection pattern is shared.

\begin{definition}[Phase object]
Let $A$ be a commutative ring with an involution $\sig$.  A phase object is a unit $U\in A$ such that
\[
    U^\sig=U^{-1}.
\]
Equivalently,
\[
    \Nm(U)=UU^\sig=1.
\]
\end{definition}

For a phase object define the unnormalized trace and spread shadows
\[
    T(U)=U+U^{-1},
\]
\[
    S(U)=U-U^{-1}.
\]
These are the native projections.  Equivalently, $T(U)$ is the trace projection of the pair $\langle U\mid U^{-1}\rangle$, while $S(U)$ is its spread projection.

\begin{proposition}[Fixed and anti-fixed projections]\label{prop:fixed-antifixed}
Let $A$ be a commutative ring with involution $\sig$, and let $U$ be a phase object.  Then
\[
    T(U)^\sig=T(U),
\]
and
\[
    S(U)^\sig=-S(U).
\]
Thus trace shadows land in the fixed part of the involution, while spread shadows land in the anti-fixed part.
\end{proposition}

\begin{proof}
Since $U^\sig=U^{-1}$, we have
\[
    T(U)^\sig=(U+U^{-1})^\sig=U^{-1}+U=T(U),
\]
and
\[
    S(U)^\sig=(U-U^{-1})^\sig=U^{-1}-U=-S(U).
\]
\end{proof}

Classical cosine and sine are later normalized versions of these quantities after choosing a complex coordinate.

\section{Chord before cosine}

Classical trigonometry often begins with cosine as a coordinate function.  The shadow-native formulation suggests a more primitive object: displacement from the identity.

Let $U$ be a norm-one phase.  Define its phase displacement
\[
    D(U)=U-1.
\]
Then the norm of this displacement is
\[
    \Nm(U-1)=(U-1)(U^{-1}-1).
\]
Expanding gives
\[
    \Nm(U-1)
    =1-U-U^{-1}+1
    =2-(U+U^{-1}).
\]
Thus:

\begin{proposition}[Chord norm identity]\label{prop:chord}
For every norm-one phase object $U$,
\[
    \Nm(U-1)=2-T(U).
\]
\end{proposition}

In a complex embedding with $U=e^{i\theta}$, this becomes
\[
    |e^{i\theta}-1|^2=2-2\cos\theta.
\]
But the shadow statement does not need $\theta$, $\cos$, or $\pi$.  It says that the trace shadow is determined by the norm of the phase displacement:
\[
    T(U)=2-\Nm(U-1).
\]

Thus trace is the complement of displacement norm.  This gives a chord-first view of trigonometry:
\begin{quote}
Trigonometry may begin with norm of displacement, not with normalized coordinate functions.
\end{quote}

\section{Trace orbits and algebraic addition formulas}

For a phase object $U$, define its trace orbit
\[
    T_n(U)=U^n+U^{-n},
\]
and spread orbit
\[
    S_n(U)=U^n-U^{-n}.
\]
When the phase object is clear, we write $T_n$ and $S_n$.

\begin{proposition}[Trace recurrence]\label{prop:trace-recurrence}
Let
\[
    T_1=U+U^{-1}.
\]
Then
\[
    T_0=2,
    \qquad
    T_{n+1}=T_1T_n-T_{n-1}
\]
for all $n\ge1$.
\end{proposition}

\begin{proof}
Multiply:
\[
    (U+U^{-1})(U^n+U^{-n})
    =U^{n+1}+U^{n-1}+U^{-(n-1)}+U^{-(n+1)}.
\]
The right-hand side is
\[
    T_{n+1}+T_{n-1}.
\]
Hence
\[
    T_1T_n=T_{n+1}+T_{n-1},
\]
which gives the recurrence.
\end{proof}

The notation $T_n$ here denotes the trace orbit, not the standard Chebyshev polynomial.  The relationship is classical:
\[
    T_n(U)=2C_n\!\left(\frac{T_1(U)}2\right),
\]
where $C_n$ is the Chebyshev polynomial of the first kind.  Equivalently, these are Dickson polynomials $D_n(T_1,1)$.  These recurrences and their trigonometric interpretations are classical; the present paper uses them as the natural trace-orbit language forced by the phase/projection viewpoint rather than claiming them as new polynomials \cite{Rivlin1990,ConwayJones1976}.

The same multiplication law gives the algebraic addition formulas.

\begin{proposition}[Shadow addition formulas]\label{prop:addition-formulas}
For integers $m,n$,
\[
    T_mT_n=T_{m+n}+T_{m-n},
\]
\[
    S_mT_n=S_{m+n}+S_{m-n},
\]
\[
    S_mS_n=T_{m+n}-T_{m-n}.
\]
\end{proposition}

\begin{proof}
The first two identities follow by expanding products of
$U^m\pm U^{-m}$ and $U^n+U^{-n}$ and grouping terms.  For the sign-sensitive identity,
\[
\begin{aligned}
S_mS_n
&=(U^m-U^{-m})(U^n-U^{-n}) \\
&=U^{m+n}-U^{m-n}-U^{-(m-n)}+U^{-(m+n)} \\
&=T_{m+n}-T_{m-n}.
\end{aligned}
\]
\end{proof}

Classical angle-addition identities are normalized versions of these shadow identities.  In the shadow layer, they come directly from the orbit law
\[
    U^mU^n=U^{m+n}.
\]

\section{Phase lift: the native replacement for arccos}

Classically, given
\[
    t=2\cos\theta,
\]
one may write
\[
    \theta=\arccos(t/2),
\]
or
\[
    e^{i\theta}=\cos\theta+i\sin\theta.
\]
Both descriptions introduce normalized coordinates and analytic inverse functions.

The shadow-native replacement is a lift relation.

\begin{definition}[Phase lift]\label{def:phase-lift}
A phase lift of a scalar $t$ is a unit $U$, equipped with the formal shadow involution $U^\sig=U^{-1}$, satisfying
\[
    U+U^{-1}=t.
\]
Algebraically, this is represented by adjoining a root of
\[
    U^2-tU+1=0.
\]
Conversely, any unit satisfying this quadratic gives such a trace relation after equipping the two roots with the involution that swaps $U$ and $U^{-1}$.  We write
\[
    \Phase(t;U)
\]
for this lift relation.
\end{definition}

Thus inverse cosine is replaced by phase lift:
\[
    \arccos(t/2)
    \quad\leadsto\quad
    U^2-tU+1=0.
\]
The lift is generally not unique: $U$ and $U^{-1}$ have the same trace.  This two-valuedness is exactly the $\sig$-orbit $\{U,U^{-1}\}$.  This is expected.  Phase is a shadow object, not a single-valued scalar function.

\section{The golden circular bridge}

We now specialize to the golden trace value.

Let $U$ be a norm-one phase satisfying
\[
    U+U^{-1}=\ph.
\]
Equivalently,
\[
    U^2-\ph U+1=0.
\]
Define
\[
    T_n=U^n+U^{-n}.
\]
By Proposition~\ref{prop:trace-recurrence},
\[
    T_{n+1}=\ph T_n-T_{n-1}.
\]

Compute:
\[
    T_0=2,
\]
\[
    T_1=\ph,
\]
\[
    T_2=\ph^2-2=(\ph+1)-2=\ph-1=\ph^{-1},
\]
\[
    T_3=\ph(\ph-1)-\ph=\ph^2-2\ph=1-\ph=-\ph^{-1},
\]
\[
    T_4=\ph(1-\ph)-(\ph-1)=1-\ph^2=-\ph,
\]
\[
    T_5=\ph(-\ph)-(1-\ph)=-\ph^2-1+\ph=-2.
\]

This gives the central theorem.

\begin{theorem}[Golden trace gives a primitive tenth phase]\label{thm:golden-pentagonal}
Let $A$ be an integral domain containing $\Z[\ph]$ as a subring, where $\ph\in A$ denotes the image of the golden ratio.  Let $U\in A^\times$ satisfy
\[
    U+U^{-1}=\ph.
\]
Then
\[
    U^4-U^3+U^2-U+1=0,
\]
so $U$ satisfies the tenth cyclotomic equation.  Consequently
\[
    U^5=-1,
    \qquad
    U^{10}=1,
\]
and $U$ has exact order $10$ in $A^\times$.
\end{theorem}

\begin{proof}
There are two useful ways to see the torsion conclusion.

\textit{Cyclotomic route.}  The golden identity and the trace equation eliminate the visible trace variable:
\begin{align*}
0&=\ph^2-\ph-1\\
  &=(U+U^{-1})^2-(U+U^{-1})-1.
\end{align*}
Multiplying this known-zero relation by $U^2$ gives
\[
    U^4-U^3+U^2-U+1=0.
\]
This polynomial is the tenth cyclotomic polynomial $\Phi_{10}(U)$.  Multiplying the cyclotomic relation by $U+1$ yields
\[
    (U+1)(U^4-U^3+U^2-U+1)=U^5+1=0,
\]
so
\[
    U^5=-1,
\]
and hence
\[
    U^{10}=1.
\]
Because $A$ contains $\Z[\ph]$ as a subring, it has characteristic zero.  Since $U^5=-1$, $U$ cannot have order $1$ or $5$.  If $U$ had order $2$, then $U^2=1$, so $(U-1)(U+1)=0$; as $A$ is an integral domain, $U=1$ or $U=-1$, giving trace $2$ or $-2$, not $\ph$.  Therefore the order is exactly $10$.

\textit{Trace-orbit route.}  The recurrence computation gives $T_5=U^5+U^{-5}=-2$.  Multiplying this equality by $U^5$ gives $(U^5+1)^2=0$, and the integral-domain hypothesis again forces $U^5=-1$.  This second route is the same torsion statement seen from the trace orbit rather than from the cyclotomic polynomial.
\end{proof}

This is pentagonal trigonometry without angles.  The algebraic equation
\[
    U+U^{-1}=\ph
\]
already forces the phase into the primitive tenth cyclotomic layer.

In the classical complex realization, one may take
\[
    U=e^{\pm i\pi/5},
\]
so that
\[
    U+U^{-1}=2\cos(\pi/5)=\ph.
\]
But in the shadow-native formulation, $\pi$, angle, sine, and cosine are not needed to obtain the torsion law.

\section{The phase-to-unit ladder}

The theorem suggests a striking ladder:
\[
    U
    \quad\longmapsto\quad
    U+U^{-1}=\ph
    \quad\longmapsto\quad
    \ph^2-\ph=1.
\]
In words:
\begin{center}
\begin{tabular}{ll}
\toprule
Hidden phase layer & $U,\ U^{-1},\ U^{10}=1$ \\
Trace projection & $U+U^{-1}=\ph$ \\
Golden scalar layer & $\ph^2-\ph=1$ \\
Rational unit & $1$ \\
\bottomrule
\end{tabular}
\end{center}

Thus the rational unit can be reached from a hidden circular phase by trace projection followed by the golden scalar identity.  This is the trigonometric analogue of the earlier projection principles:
\[
    \text{addition} = \text{trace projection},
\]
\[
    \text{quotient derivative} = \text{factor projection},
\]
\[
    \text{golden scalar} = \text{phase trace projection}.
\]

\section{Pentagonal trigonometry without angles}

The trace orbit of a golden phase closes algebraically:
\[
    T_0=2,
\]
\[
    T_1=\ph,
\]
\[
    T_2=\ph^{-1},
\]
\[
    T_3=-\ph^{-1},
\]
\[
    T_4=-\ph,
\]
\[
    T_5=-2.
\]
Then
\[
    T_{n+10}=T_n.
\]
Classically this is the list
\[
    2\cos(n\pi/5),\qquad n=0,1,2,3,4,5.
\]
But the native construction uses only trace recurrence and the golden identity.

This is the basic pentagonal trigonometric table:
\begin{center}
\begin{tabular}{c c c}
\toprule
$n$ & $T_n=U^n+U^{-n}$ & classical interpretation \\
\midrule
0 & $2$ & $2\cos 0$ \\
1 & $\ph$ & $2\cos(\pi/5)$ \\
2 & $\ph^{-1}$ & $2\cos(2\pi/5)$ \\
3 & $-\ph^{-1}$ & $2\cos(3\pi/5)$ \\
4 & $-\ph$ & $2\cos(4\pi/5)$ \\
5 & $-2$ & $2\cos \pi$ \\
\bottomrule
\end{tabular}
\end{center}

The right column is an interpretation, not a definition.  Algebraically, these are the visible traces of powers of the phase lift.

\section{Spread shadows and sine-like data}

Define
\[
    S_n=U^n-U^{-n}.
\]
Then
\[
    S_n^2=(U^n-U^{-n})^2
\]
\[
    =U^{2n}-2+U^{-2n},
\]
whereas
\[
    T_n^2=(U^n+U^{-n})^2
\]
\[
    =U^{2n}+2+U^{-2n}.
\]
Therefore
\[
    S_n^2=T_n^2-4.
\]

\begin{proposition}[Spread-square identity]\label{prop:spread-square}
For every norm-one phase object $U$,
\[
    S_n^2=T_n^2-4.
\]
\end{proposition}

In the complex circular interpretation,
\[
    S_n=2i\sin(n\theta).
\]
But the shadow-native object is simply the spread projection.  Normalized sine requires dividing by $2i$; this is a fixed-denominator rate/localization step followed by a coordinate choice, not a variable reciprocal operation.

Thus sine-like information is not absent.  It is present as unnormalized spread.

\section{Tangent as relation, not division}

Classically,
\[
    \tan\theta=\frac{\sin\theta}{\cos\theta}.
\]
In trace/spread variables this is a quotient between spread and trace.  To stay consistent with the quotient discipline of the second paper, we should not make tangent primitive.

Instead, define a tangent parameter relationally.  After choosing a complex embedding, or after adjoining a skew unit $i$ satisfying $i^\sig=-i$ and $i^2=-1$, one may write
\[
    S(U)=i\tau T(U).
\]
The parameter $\tau$ plays the role of tangent when the relation has a unique solution.  Uniqueness requires a cancellation condition on $T(U)$, just as quotient relations require nonzero-divisor hypotheses in general rings.

The point is methodological:
\begin{quote}
Tangent is a slope relation between spread and trace, not a primitive division operation.
\end{quote}

This also clarifies where tangent singularities come from.  They arise when the trace projection vanishes, not because tangent was a fundamental object.

\section{The double life of \texorpdfstring{$\varphi$}{phi}}

The golden ratio now has two complementary roles.

\subsection{The hyperbolic life: real quadratic growth}

In the real quadratic setting, the golden pair is
\[
    \Phi=\langle \ph\mid\ps\rangle,
    \qquad
    \ps=1-\ph=-\ph^{-1}.
\]
Its powers yield Lucas and Fibonacci data:
\[
    \Tr(\Phi^n)=\ph^n+\ps^n=L_n,
\]
\[
    \Sp(\Phi^n)=\ph^n-\ps^n=\sqrt5\,F_n.
\]
This is a hyperbolic or growth-like role.  The unit $\ph$ generates expansion and contraction along real embeddings.

\subsection{The circular life: pentagonal phase trace}

In the circular setting, $\ph$ is no longer the phase.  It is the visible trace of a hidden phase:
\[
    U+U^{-1}=\ph.
\]
Then Theorem~\ref{thm:golden-pentagonal} gives
\[
    U^5=-1,
    \qquad
    U^{10}=1,
\]
with exact order $10$ in the integral-domain setting.  This is the pentagonal or circular role.

Thus the golden ratio lives a double life:
\begin{quote}
It is a generator of real quadratic growth, and also the footprint of discrete circular rotation.
\end{quote}

This duality is one of the main conceptual payoffs of the shadow framework.

\section{Cyclotomic traces and the classical interpretation}

Let $\zeta_{10}$ be a primitive tenth root of unity.  In the complex embedding one may choose
\[
    \zeta_{10}=e^{i\pi/5}.
\]
Then
\[
    \zeta_{10}+\zeta_{10}^{-1}=2\cos(\pi/5)=\ph.
\]
Similarly, for a primitive fifth root $\zeta_5=e^{2\pi i/5}$,
\[
    \zeta_5+\zeta_5^{-1}=2\cos(2\pi/5)=\ph-1=\ph^{-1}.
\]

These are standard cyclotomic identities: real subfields of cyclotomic fields are generated by sums of the form
\[
    \zeta+\zeta^{-1}.
\]
This is classical cyclotomic-field theory \cite{Washington1997}.  The present paper does not claim these identities as new.  The new point is their role in the $\ph$-native hierarchy: the golden scalar appears as a trace projection from a circular phase object.

\section{Delaying \texorpdfstring{$\pi$}{pi} and continuous phase}

In the algebraic layer, period is torsion:
\[
    U^{10}=1.
\]
There is no need to introduce $\pi$ to state or prove this relation.

After choosing a complex embedding and a continuous phase coordinate, one writes
\[
    U=e^{i\pi/5}.
\]
At that point $\pi$ enters as a coordinate of continuous phase, not as an algebraic primitive.

There are two distinct completions here.

\subsection{Torsion completion}

Under the standard complex embedding, the collection of all roots of unity is dense in the unit circle.  Thus continuous circular phase can be approached topologically by allowing torsion phases of all orders.  This is not a claim that the single golden/pentagonal tower is dense by itself; the golden phase supplies one distinguished torsion orbit of exact order $10$.  In this broader torsion completion, angle emerges as a completion of finite cyclic phase data.

\subsection{Differential or flow completion}

Topological density of torsion phases is not yet calculus.  To obtain smooth rotation, angular velocity, differential equations, and Fourier analysis, one also needs a continuous flow structure.  This may be where the factor-derivative machinery from \emph{Quotients Without Division} re-enters, now applied to norm-one phase orbits.

Thus the tentative answer is:
\[
\boxed{\text{continuous phase} = \text{torsion completion} + \text{differential flow structure}.}
\]

This paper develops the algebraic and torsion side.  The analytic flow side is left for later work.

\section{Comparison with classical trigonometry}

The translation is now clear.

\begin{center}
\begin{tabular}{p{0.36\linewidth}p{0.45\linewidth}}
\toprule
Classical concept & Shadow-native concept \\
\midrule
angle $\theta$ & phase object $U$ or orbit index \\
period $2\pi$ & torsion relation $U^n=1$ \\
$e^{i\theta}$ & norm-one phase \\
$2\cos\theta$ & trace $U+U^{-1}$ \\
$2i\sin\theta$ & spread $U-U^{-1}$ \\
$\arccos(t/2)$ & phase lift $U^2-tU+1=0$ \\
tangent & spread--trace slope relation \\
addition formulas & multiplication of phase units \\
Chebyshev polynomials & trace recurrences \\
$\pi$ & analytic coordinate of phase completion \\
\bottomrule
\end{tabular}
\end{center}

The native layer is unnormalized and algebraic.  The classical layer is normalized and analytic.  Both are useful, but they should not be confused.

\section{Open questions}

\begin{question}[General torsion traces]
Classically, if $\zeta$ is a primitive $n$-th root of unity, then
\[
    \Q(\zeta+\zeta^{-1})
\]
is the maximal real subfield of the $n$-th cyclotomic field, typically of degree $\phi_{\mathrm E}(n)/2$, where $\phi_{\mathrm E}$ denotes Euler's totient function.  Thus the question of when $\zeta+\zeta^{-1}$ lies in $\Q(\sqrt5)$, and when it lies in the integer ring $\Z[\ph]$, is a tractable cyclotomic-field classification problem.  What does the shadow grammar viewpoint recover about this classification, and can it distinguish the golden cases conceptually rather than only computationally?
\end{question}

\begin{question}[Native sine normalization]
Can spread shadows
\[
    S(U)=U-U^{-1}
\]
be used directly in applications, or is normalized sine unavoidable once analytic geometry enters?
\end{question}

\begin{question}[Phase-lift algebra]
What is the correct type-theoretic or categorical status of the phase-lift relation
\[
    U^2-tU+1=0?
\]
Is it a relation, a dependent pair, or a new phase sort?
\end{question}

\begin{question}[Tangent without division]
After choosing a complex embedding or adjoining a skew unit $i$, can tangent be developed usefully as a relation
\[
    S(U)=i\tau T(U)
\]
without introducing quotient singularities as primitives?
\end{question}

\begin{question}[Continuous phase]
Should continuous angle be constructed as a limit of high-order torsion orbits, through an explicit differential flow relation, or through both?
\end{question}

\begin{question}[Golden Fourier theory]
If $\ph$ is the trace of pentagonal phase, is there a useful $\ph$-native Fourier theory based on finite phase orbits and trace/spread projections?
\end{question}

\section{Conclusion: the golden ratio as circular trace}

Classically one writes
\[
    2\cos(\pi/5)=\ph.
\]
The shadow-native formulation reverses the emphasis.  The golden ratio is a visible scalar, and the hidden phase $U$ satisfying
\[
    U+U^{-1}=\ph
\]
is its circular lift.  This lift obeys
\[
    U^5=-1,
    \qquad
    U^{10}=1.
\]

Thus $\ph$ is not only a real quadratic unit generating growth.  It is also the trace projection of pentagonal circular phase.

This completes the third step of the program.  In the first paper, addition became trace projection.  In the second, quotient-like operations split into unit action, relation, approximation, normalization, and factor projection.  In the present paper, trigonometry becomes projection theory of norm-one phase orbits.

The final message is:
\begin{quote}
Trigonometry begins not with angles, but with phase shadows whose traces become visible arithmetic.
\end{quote}

\begin{thebibliography}{9}

\bibitem{SantosGoldenShadow}
V.~Santos.
\newblock \emph{A Golden Shadow Algebra: From a $\varphi$-Leaf Scalar Core to Conjugate-Pair Arithmetic}.
\newblock Exploratory manuscript, 2026.

\bibitem{SantosQuotients}
V.~Santos.
\newblock \emph{Quotients Without Division: Unit Actions, Relations, and Factor Projections in a $\varphi$-Native Algebra}.
\newblock Exploratory manuscript, 2026.

\bibitem{Odrzywolek2026EML}
A.~Odrzywo\l{}ek.
\newblock \emph{All elementary functions from a single binary operator}.
\newblock arXiv:2603.21852, 2026.

\bibitem{Washington1997}
L.~C. Washington.
\newblock \emph{Introduction to Cyclotomic Fields}.
\newblock 2nd edition, Springer, 1997.

\bibitem{Rivlin1990}
T.~J. Rivlin.
\newblock \emph{Chebyshev Polynomials: From Approximation Theory to Algebra and Number Theory}.
\newblock 2nd edition, Wiley, 1990.

\bibitem{ConwayJones1976}
J.~H. Conway and A.~J. Jones.
\newblock Trigonometric Diophantine equations.
\newblock \emph{Acta Arithmetica}, 30:229--240, 1976.

\end{thebibliography}

\end{document}
